\chapter{Metodologia Prática}

\section{Aula 1}
O primeiro curso de \LaTeX  da UDESC CEPLAN ministrado pelo professor Lucas Zick no dia \today foi um sucesso, nesta aula a gente aprendeu sobre a história do \LaTeX, que se iniciou com a uma frustração do professor Donald Knuth vendo as provas impressas dos seus livros, Donald Knuth criou então o TeX que era uma linguagem de diagramação, porém a maior dificuldade e empecilho do TeX é que essa linguagem era muito pesada para época, era muito difícil de aprender e de usar e necessariamente precisava ser instalada na máquina. Vendo isso, Leslie Lamport resolveu criar uma camada de alto nível na linguagem TeX, facilitando demais o uso dessa linguagem e tendo uma curva de aprendizado mais fácil, criando assim o \LaTeX que conhecemos hoje (Lamport TeX).

Mesmo com uma curva de aprendizado mais fácil o \LaTeX ainda não é muito difundido pois ele ainda sim impõem medo nas pessoas no primeiro contato delas com a linguagem, muitas pessoas não entendem o porque você tem que aprender uma linguagem de diagramação sendo que criar um documento word parece muito mais intuitivo, porém a padronização do \LaTeX dentro de um documento muito extenso como por exemplo um TCC é absurdamente mais fácil, o documento não quebra ao inserir uma imagem nova, as referências e o sumário por exemplo é algo que você não tem que se preocupar pois é automático, etc...

Nessa aula nós iniciamos o aprendizado da linguagem de diagramação dentro do site OverLeaf, aprendemos a programar o preâmbulo, também algumas funções e comandos básicos que podemos usar no \LaTeX com o uso de barra invertida como a quebra de linha\\que podemos ver nessa frase, pois o \LaTeX não quebra a linha com um enter. Outros comandos que a gente também aprendeu foi o \textbf{texto em negrito}, \textit{texto em itálico}, \underline{texto sublinhado}, aprendemos também a colocar alguns caracteres especiais que tem uma forma mais "chatinha" do que só digitar ele como o \%, \$, \&, entre outros caracteres especiais.


\section{Aula 2}
Nesta segunda aula nós tivemos muito mais conteúdos práticos, basicamente o objetivo da aula é sair de um documento básico e padrão, o famoso "hello word" de uma única página para uma estrutura muito mais completa. Para isso nós aprendemos os modelos/estilos padrões de documentos que possuímos no \LaTeX, como o artigo curto (article) que são os artigos mais simples, sem folha de rosto e sem capítulos por exemplo, artigo longos (report) que são documentos mais completos, geralmente utilizados em TCCs e Dissertações, eles já possuem folha de rosto, sumários e capítulos, esse modelo de documento é justamente o modelo que eu estou utilizando para a escrita desse meta-artigo do curso, e por último temos o livro (book) que é o modelo mais "completo" e é o modelo que exige mais conteúdo de todos eles, além disso o livro possui algumas coisas que os outros modelos não possuem como uma otimização para ser impresso em folhas frente e verso, possuindo margens assimétricas e forçando que os capítulos se iniciem em páginas ímpares.

Depois da criação do modelo nós aprendemos alguns novos parâmetros para utilizar no preâmbulo, logo na definição do modelo é possível definir por exemplo tamanho da fonte (Este documento está definido em 12pt), estilo da folha como folha A4 (a4paper) e folha para cartas (letterpaper) (este documento está definido como folha A4). Aprendemos também a definir as margens da folha de acordo com as normas ABNT e a adicionar afiliações nos autores do documento, como por exemplo neste documento eu como autor estou afiliado com a UDESC, e adicionei meu email intitucional que ele fica no rodapé da capa.

Outras coisas que aprendemos nessa aula é a adição de resumo e sumário, esses dois comandos mostram que o \LaTeX realmente ajuda muito na escrita de documentos complexos como um TCC, utilizando apenas um comando de \textbackslash abstract nós podemos escrever texto dentro dessa seção que toda a formatação correta do resumo fica pronta automaticamente, e um \textbackslash tableofcontents que o sumário magicamente aparece completo e sem erro automaticamente.

Agora falando das seções de conteúdos nós aprendemos a criar capítulos com o comando \textbackslash chapter, seções com o comando \textbackslash section, subseções com o comando \textbackslash subsection e por fim sub sub seções com o comando \textbackslash subsubsection que funcionam exatamente nessa hierarquia, é importante saber que nos artigos curtos não existe os capítulos e nos artigos longos e livros as sub sub seções não existem também. Por fim um detalhe importante que podemos fazer é excluir a numeração das seções utilizando um * no final do comando (\textbackslash section*\{\}), tirando a numeração a seção desaparece do sumário, porém podemos adicionar ela denovo com esse comando (\textbackslash addcontentsline\{toc\}\{section\}\{Introdução\})

\section{Aula 3}
Esta aula foi recheada de novas funções e comandos, dentre essas funções temos a alteração de fontes das letras, diferentes tipos de enfatização nos fluxos de textos, diferentes variações de tamanhos de textos, diferentes alinhamentos de texto, diferentes espaçamentos entre as linhas e espaçamento no início dos parágrafo além de espaçamentos verticais e horizontais de forma forçada, formatação de citações longas e por fim o uso de diferentes listas.

Começando com as diferentes Famílias de fontes nativas do \LaTeX, o \LaTeX tem por padrão um formato de letras que são serifadas, isso quer dizer que as letras no documento vão possuir pequenos "pezinhos" ou cones nas extremidades, o que deixa as letras um pouco mais elegantes mas que era principalmente muito importante na época das máquinas de datilografia porque as extremidades das letras concentravam mais quantidade de tinta e mesmo que a letra saísse mais apagada as pessoas conseguiam ler as letras mais facilmente por conta das extremidades possuirem mais tinta. Outra família de fonte que o \LaTeX possui nativamente é o texto monoespaçado, este modelo significa que todas as letras vão ocupar o mesmo tamanho, esse estilo de texto é muito usado em trechos de códigos e URLs pois as letras ficam mais espaçadas e fica mais fácil de identifica-las. Aqui temos exemplos desses diferentes tipos de texto:\\

\noindent Este é um texto com serifas nas bordas.\\
\textsf{Este é um texto sem serifa nas bordas.}\\
\texttt{Este é um texto monoespaçada ótimo para destacar trechos de códigos ou URLs.}

Outras funções que possuímos nativamente são os diferentes tipos de enfatização, as enfatizações compões o negrito (\textbackslash textbf\{\}) que deixa o texto mais forte e mais grosso ficando mais visível dentro do texto, o itálico \textbackslash textit\{\}) que deixa o texto mais "deitado" que é muito utilizado para destacar palavras com sentidos especiais ou palavras estrangeiras, temos também o sublinhado (\textbackslash underline\{\}) que adiciona uma linha por baixo do texto, e por fim temos o Small Caps (\textbackslash textsc\{\}) que deixa o texto em Caps Lock porém pequeno, esses são os exemplos:\\

\noindent \textbf{Este é um texto em negrito.}\\
\textit{Este é um texto em itálico.}\\
\underline{Este é um texto sublinhado.}\\
\textsc{Este é um texto em Small Caps.}\\
ESTE É UM TEXTO EM CAPS LOCK PARA COMPARAÇÃO.\\

O \LaTeX possui também diferentes variações de tamahos de texto, esses tamanhos são relativos e variam de acordo com o tamanho base do texto que foi definido no início do documento, neste artigo o tamanho base do documento está definido como 12pt e o tamanho dos textos vão variar de forma relativa a esse tamanho. Os comandos são \textbackslash tiny, \textbackslash scriptsize, \textbackslash small, \textbackslash normalsize, \textbackslash large, \textbackslash Large, \textbackslash huge, \textbackslash Huge e estes são exemplos de todos os tamanhos:\\

\noindent {\tiny Este é um texto utilizando o tamanho minúsculo.}\\
{\scriptsize Este é um texto utilizando o tamanho de script.}\\
{\small Este é um texto utilizando o tamanho pequeno.}\\
{\normalsize Este é um texto utilizando o tamanho normal padrão.}\\
{\large Este é um texto utilizando o tamanho grande.}\\
{\Large Este é um texto utilizando o tamanho muito grande.}\\
{\huge Este é um texto utilizando o tamanho gigante.}\\
{\Huge Este é um texto utilizando o tamanho colossal.}\\

Nós podemos também utilizar diferentes tipos de alinhamentos no texto, por padrão o \LaTeX utiliza o alinhamento justificado para o textos, que basicamente é o modelo onde as linhas comecam e terminam sempre no mesmo lugar, deixando o texto mais quadrado, mas podemos utilizar diferentes alinhamentos nativamente. Para modificar o alinhamento dos parágrafos nós fazemos o uso de ambientes, onde nós vamos fazer a abertura desse ambiente utilizando o comando \textbackslash begin\{\} e obrigatoriamente o fechamento desse ambiente utilizando o comando \textbackslash end\{\}, dentro dessas chaves nós colocamos a variável que vai ditar o alinhamento do ambiente, podendo ser "center", "flushleft" e "flushright", aqui temos exemplos desses alinhamentos:\\

\begin{center}
    Este é um texto com alinhamento centralizado.
\end{center}

\begin{flushleft}
    Este é um texto com alinhamento à esquerda com irregularidade à direita.
\end{flushleft}

\begin{flushright}
    Este é um texto com alinhamento à direita com irregularidade à esquerda.
\end{flushright}

Outro detalhe importante que podemos alterar no documento é o espaçamento entre as linhas, por padrão o latex utiliza um espaçamento de linhas de 1.0, porém nós podemos alterar esse espaçamento importando um pacote direto no preâmbulo, esse pacote se chama \textbackslash usepackage\{setspace\}, depois da importação do pacote nós podemos configurar o espaçamento utilizando os seguintes comandos logo depois da importação do pacote:\\

\noindent \textbackslash onehalfspacing para espaçamento de 1.5 (padrão ABNT).\\
\textbackslash doublespacing para espaçamento duplo.\\
\textbackslash singlespacing para espaçamento de 1.0 (padrão do \LaTeX).\\

Além do espaçamento de linhas nós podemos configurar o espaçamento de recuo que é o recuo que vem na primeira linha de cada parágrafo, o \LaTeX já possui um recuo padrão mas podemos configurar esse recuo com o comando \textbackslash setlength\{\textbackslash parindent\}\{1.25cm\} por exemplo. E se caso não quisermos um recuo podemos iniciar a frase com o comando \textbackslash noindent.

Ainda falando de recuo e espaçamento as vezes precisamos forçar um espaçamento empurrando algum texto para baixo ou para os lados, por exemplo pra criar a folha de rosto. Para isso nós podemos utilizar alguns comandos como o \textbackslash vspace\{\} e \textbackslash hspace\{\}:

Primeiro texto

\vspace{2cm}
Este é o segundo texto depois do espaçamento de 2cm

Assinatura: \hspace{5cm} Data:

No \LaTeX, a nota de rodapé é declarada exatamente onde a marcação deve aparecer no texto. O compilador joga o texto lá para baixo sozinho, exemplo:

A Universidade do Estado de Santa Catarina \footnote{Fundada em 1965, é uma das principais intituições estaduais do país.} possui diversos campos espalhados pelo estado.

Outra ferrameta muito interessante do \LaTeX é a formatação automática de citações longas. A ABNT exige um recuo de 4cm, fonte menor (10pt) e espaçamento simples quando alguma citação direta passa de 3 linhas, e para fazer essa formatação de forma rápida e simples utilizando o comando \textbackslash begin\{quote\} e dentro do ambiente de citação nó adicionamos dois comandos que é o \textbackslash singlespacing e o \textbackslash small:
\begin{quote}
    % \singlespacing 
    \small
    "E aqui está a formatação e o uso desse estilo de citação direta que passa de 3 linhas, esse tipo de formatação se fosse feita à mão daria um trabalho e isso já está pronto utilizando apenas alguns comandos nossa citação já está pronta."
\end{quote}

Nós podemos também criar listas dentro do \LaTeX, isso é criado através de ambientes, um exemplo de listas é listas usando bullet points, que são as listas não numeradas que ficam apenas com umas bolinhas representando os itens da lista, essas listas são criadas com ambientes de itemize:

\begin{itemize}
    \item item 1
    \item item 2
\end{itemize}

Mas podemos criar listas numeradas também, para listas numeradas usamos ambiente de enumerate:

\begin{enumerate}
    \item item 1
    \item item 2
\end{enumerate}

Podemos também criar listas dentro de listas, misturando os modelos das listas como nesse exemplo:

\begin{itemize}
    \item Item 1 da primeira lista
    \begin{enumerate}
        \item Item 1 da segunda lista
        \item Item 2 da segunda lista
    \end{enumerate}
    \item Item 2 da primeira lista
\end{itemize}

Importando o pacote enumitem nós podemos modificar a numeração das listas numeradas, podendo ser no modelo alfabetico (\textbackslash alph*) e maiúsculo (\textbackslash Alph*), e no modelo romano minúsculo (\textbackslash roman*) e maiúsculo (\textbackslash Roman*) 

\section{Aula 4}
A aula 4 foi focada na escrita e criação de fórmulas matemáticas dentro do \LaTeX. Primeiramente, para podermos usufruir do potêncial total das fórmulas matemáticas nós precisamos importar 3 dos principais pacotes de matemática, sendo eles \textbackslash usepackage\{amsfonts\}, \textbackslash usepackage\{amsmath\} e \textbackslash usepackage\{amssymbol\}.

Existem dois modos de criação de equações no \LaTeX, o modo inline e modo display. No modo inline a equação ela é introduzida dentro do próprio fluxo de texto, sem quebrar linha e sem nenhum espaçamento entre as linhas, para adicionarmos uma equação ou fórmula matemática nós precisamos criar um ambiente matemático que é iniciado com um cifrao (\$) e finalizado com outro cifrão, como por exemplo:

Se escrevermos a frase "Calculamos a area do circulo como \$A = \textbackslash pi r\^{}2\$." nó teremos como resultado "Calculamos a area do circulo como $A = \pi r^2$."

O outro modo de adicionarmos funções e equações matemáticas no nosso texto é o modo display, onde diferente do modo inline o  texto é interrompido com uma quebra de linha e um espaçamento e a fórmula é centralizada na nova linha, ela também usa tamanhos maiores. Para adicionarmos equações no modo display nós iniciamos o ambiente utilizando barra e colchetes (\textbackslash[) e finalizamos com colchetes barra (]\textbackslash), como por exemplo:

Se escrevermos a frase "A formula de Bhaskara e expressa abaixo: 2 \textbackslash [x = \textbackslash frac{-b \textbackslash pm \textbackslash sqrt{b\^{}2- 4ac}}{2a}\textbackslash ] 3 Logo, podemos encontrar as raizes." nós teremos como resultado "A formula de Bhaskara e expressa abaixo: 2 \[x = \frac{-b \pm \sqrt{b^2- 4ac}}{2a}\] 3 Logo, podemos encontrar as raizes."

Para adicionarmos textos sobrescritos, que é o número que representa a potência e vai ficar no canto superios de uma variável ou número, nós usamos o acento "chapéu" (\^{}). Um exemplo disso, ao usarmos a frase "O Teorema de Pitagoras e dado por \$c\^{}2 = a\^{}2 + b\^{}2\$." teremos como resultado "O Teorema de Pitagoras e dado por $c^2 = a^2 + b^2$.".

Além do sobrescrito nós temos o subscrito, que é o número que representa o índice de uma variável e vai ficar no canto inferior de uma variável ou número, para utilizarmos o subscrito nós usamos o underline (\_), como por exemplo se escrevermos isso "O robo iniciou na posicao \$x\_0\$ e foi ate \$x\_1\$." teremos como resultado "O robo iniciou na posicao $x_0$ e foi ate $x_1$."

É muito importante durante o uso dessas duas funções nós adicionarmos chaves externas para representar mais de um algarismo no texto subscrito ou no texto sobrescrito porque se não o \LaTeX não vai subentender por padrão que apenas o primeiro algarismo vai ser representado na potência ou no índice. Por exemplo: se escrevermos \$a\^{}22\$ o resultado é = $a^22$, já se escrevermos \$a\^{}\{22\}\$ o resultado é = $a^{22}$.

Para adicionar frações nós utilizamos o comando \textbackslash frac\{numerador\}\{denominador\} onde a parte de cima da barra da fração é o numerador e a parte de baixo é o denominador. Então por exemplo se escrevermos \$\textbackslash frac\{1\}\{2\}\$ nosso resultado vai ser $\frac{1}{2}$

Para adicionarmos raízes nós usamos o comando \textbackslash sqrt[opcional]\{conteudo\}, onde dentro dos colchetes o valor opcional que podemos adicionar representa o coeficiente da fração, logo se não adicionarmos nenhum valor o coeficiente padrão que o \LaTeX utiliza é o 2. Por exemplo se escrevermos "\$\textbackslash sqrt\{25\} = 5\$" nós temos como resultado "$\sqrt{25} = 5$".

Na matemática é muito comun nós utilizarmos letras do alfabeto grego, por exemplo pra representar ângulos ou pra representar o "pi". Essas letras gregas nós podemos adicionar escrevendo elas por extenso logo depois de uma barra invertida, por exemplo:

\begin{itemize}
    \item \textbackslash alpha = $\alpha$
    \item \textbackslash beta = $\beta$
    \item \textbackslash gamma = $\gamma$
    \item \textbackslash theta = $\theta$
    \item \textbackslash pi = $\pi$
    \item \textbackslash omega = $\omega$
\end{itemize}

Além das letras gregas nós podemos adicionar símbolos que representam somatório (\textbackslash sum) e produtório (\textbackslash prod), e nesses símbolos em específicos o uso de texto subscrito e sobrescrito possuem um comportamento um pouco diferente do que em variáveis normais, onde o texto sobrescrito vai ficar literalmente em cima do símbolo e o texto subscrito vai ficar literalmente em baixo do símbolo. Por exemplo nesta função "\textbackslash[\textbackslash sum\_\{i=1\}\^{}\{n\} i\^{}2\textbackslash] o resultado vai ser:

\[\sum_{i=1}^{n} i^2\] 
%\[\prod_{i=1}^{n} (x_i- \bar{x})\]

Além de somatório e produtório nós podemos adicionar limites e integrais com o uso dos comandos \textbackslash lim e \textbackslash int. Por exemplo na função "\textbackslash[\textbackslash int\_\{a\}\^{}\{b\} x\^{}2 \textbackslash , dx\textbackslash ]" o resultado seria:

\[\int_{a}^{b} x^2 \, dx\]
%\[\lim_{x \to \infty} \frac{1}{x} = 0\]

Além de números e variáveis dentro das equações nós podemos também utilizar textos. Se nós somente escrevermos o texto normal dentro das equações o \LaTeX vai achar que são variáveis deixando a fonte em itálico e com formatação de variáveis, o que fica esquisito, para isso não acontecer nós utilizamos o comando \textbackslash text\{\}. Por exemplo na função "\textbackslash [\textbackslash text\{Velocidade\} = \textbackslash frac\{\textbackslash text\{Distancia\}\}\{\textbackslash text\{Tempo\}\}\textbackslash ]" o resultado seria:

\[\text{Velocidade} = \frac{\text{Distancia}}{\text{Tempo}}\]

Em artigos científicos e acadêmicos é muito comum nós enumerarmos nossas equações igual fazemos com figuras e tabelas por exemplo, para que durante o texto nós possamos referenciá-las mais facilmente. Para isso nós em vez de criar um ambiente de equação com \textbackslash[]\textbackslash nós utilizamos um \textbackslash begin\{equation\} e para que possamos referenciar depois nós adicionamos um \textbackslash label\{\} no final do ambiente passando como parâmetro dentro da chave um nome que será utilizado para referenciá-lo, assim todas as equações no texto estarão numeradas na margem direita de forma automática. Este é um exemplo de uma equação numerada e referenciada utilizando como nome o \textbackslash label\{eq:einstein\}:

\begin{equation}
    E = mc^2
    \label{eq:einstein}
\end{equation}

Como visto na Equacao \ref{eq:einstein} (foi utilizado o \textbackslash ref\{eq:einstein\}), a energia...

Outra coisa que o \LaTeX faz que salva vida é quando precisamos desenhar uma dedução matemática que ocupa várias linhas onde todas as linhas tem que ser alinhadas pelo sinal de igual (=), e isso nós conseguimos facil muito facilmente utilizando um ambiente \textbackslash begin\{align\} onde nós utlizamos o símbolo \& para marcar o ponto de alinhamento e \textbackslash\textbackslash para pular linha. Por exemplo:

\begin{align}
    2x - 5y &= 8 \\
    3x + 9y - 99z &= -12
\end{align}

Por fim nós podemos criar matrizes que são estruturadas como tabelas matemáticas. As matrizes também são muito simples de se criar e também são criadas utilizando ambiente \textbackslash begin\{matrix\} dentro de um ambiente de equação (\textbackslash[]\textbackslash)

\[\begin{bmatrix} 1 & 2 & 3 \\ 4 & 5 & 6 \\ 7 & 8 & 9 \end{bmatrix}\]


No final da aula o professor propôs dois desafios, o primeiro desafio foi representar uma matriz de linhas de tamanho m e colunas de tamanho n:

\[\begin{pmatrix}
    a_{11} & a_{12} & \cdots & a_{1n}\\
    a_{21} & a_{22} & \cdots & a_{2n}\\
    \vdots & \vdots & \ddots & \vdots\\
    a_{m1} & a_{m2} & \cdots & a_{mn}
\end{pmatrix}\]

E o segundo desafio foi escrever a equação de Schrödinger que representa o diferencial parcial que governa a função de onda de um sistema mecânico-quântico não relativístico:

\[
    i\hbar \frac{\partial}{\partial_t} \Psi (\mathbf{r}, t) = \left[ \frac{-\hbar^2}{2m} \nabla^2 + V(\mathbf{r}, t) \right] \psi (\mathbf{r}, t)
\]

\section{Aula 5}

Na quinta aula do nosso curso de \LaTeX nós aprendemos sobre inserção de imagens, figuras e tabelas dentro do nosso documento. Na edição de documentos como o Word as pessoas sofrem vários problemas quando o assundo é a inserção de imagens, muitas vezes as imagens não cabem e acabam pulando as páginas deixando um grande espaço em branco na página anterior, dependendo de como adicionamos imagens ou parágrafos a formatação acaba quebrando, e em um documento com várias imagens que estão todas já com numeração, se você precisar colocar uma imagem a mais no meio do texto você vai precisar corrigir a numeração uma por uma manualmente, e todos esses problemas são corrigidos no \LaTeX automaticamente.

No \LaTeX as imagens e figuras por padrão são consideradas como objetos flutuantes, e isso elimina aquele principal problema que eu comentei, se uma imagem não cabe no fim da página 3, o LaTeX continua puxando o texto que viria depois dela para preencher a página 3, e posiciona a imagem de forma elegante no topo da página 4, eliminando todos os espaços em brancos causados pelas imagens. Essa flutuação das imagens pode ser desativadas importando o pacote float (\textbackslash usepackage\{float\}).

Para adicionarmos figuras e imagens, primeiramente nós precisamos importar um pacote chamado graphicx, esse pacote é essencial para que o \LaTeX consiga importar arquivos externos de imagens, como PNG, JPG, PDF, etc.. Depois disso nós já podemos baixar essas imagens dentro do projeto e importar elas para o documento utilizando a função \textbackslash includegraphics[opcoes]\{caminho/imagem\}, sendo que nessas [opcoes] nós podemos mexer no tamanho da imagem, tando definindo tamanho em pixels([width=XX]) quanto definindo um tamanho em relação ao tamanho original da imagem ([scale=XX]), como por exemplo:

\includegraphics[scale=0.15]{figures/marca-udesc.jpg}

Porém dentro de um documento nós não apenas largamos uma imagem no meio do texto, nós precisamos seguir padrões e principalmente colocar uma legenda e uma numeração para cada imagem dentro do documento, e para isso nós utilizamos um ambiente de figuras (\textbackslash begin\{figure\}). Dentro desse ambiente nós podemos utilizar de algumas funções como a função \textbackslash centering que centraliza a imagem, \textbackslash includegraphics[opcional]\{caminho/imagem\} que é para adicionar a imagem, \textbackslash caption\{Texto de descrição da imagem\} que é para colocar uma pequena descrição/título para a imagem e por fim \textbackslash label\{fig:apelido\} que adiciona um apelido para a imagem para ficar mais fácil pra referencia-la durante o desenvolvimento do documento.
% \includegraphics[scale=1]{figures/zickinho.jpg}
% \includegraphics[width=0.4\textwidth]{figures/zickinho.jpg}

\begin{figure}[htbp]
    \centering
    \includegraphics[width=0.5\linewidth]{figures/zick.jpg}
    \caption{Descrição de exemplo}
    \label{fig:zickinho}
\end{figure}

A figura~\ref{fig:zickinho} é uma imagem que eu peguei no insta do nosso professor que está lecionando este curso de \LaTeX.

Por fim existe uns valores opcionais que nós podemos adicionar nas imagens e futuramente nas tabelas que apesar de serem opcionais eles são bem recomendados, estes valores são "htbp". Este valor representa a ordem de alocação da imagem na página, o "h" representa "here" que quer dizer que primeiramente o \LaTeX vai ter que tentar colocar a imagem no lugar do texto que ela foi adicionada, o "t" representa "top", ou seja, se a imagem não pode ser alocada no lugar que ela foi adicionada o \LaTeX tem que tentar colocar essa imagem no topo da página, mesma coisa para o "b" de "bottom", e por fim o "p" é de "page", então se não for possível colocar em nenhum desses lugares informados anteriormente a imagem vai ser colocada em uma página exclusiva para ela.

Nesta aula também aprendemos a inserir tabelas, as tabelas tem uma ambientação diferente pois diferente de imagens nós utilizamos dois ambientes um dentro do outro. O primeiro ambiente é o ambiente table (\textbackslash begin\{table\}) responsável por fazer a tabela flutuar na página, centralizar e receber o \textbackslash caption e o \textbackslash label. Já o outro ambiente é responsável pela formação da grade das colunas e linhas da tabela, e para isso nós utilizamos o ambiente tabular (\textbackslash begin\{tabular\}).

Então primeiramente a gente abre o ambiente table onde já podemos utilizar os valores opcionais de alocação da imagem (\textbackslash begin\{table\}[htbs]), depois nós utilizamos a mesma função das imagens de centralização que é o \textbackslash centering, depois nós já podemos desenhar a tabela com o ambiente tabular, para isso nós iniciamos abrindo o ambiente tabular \textbackslash begin\{tabular\} onde este ambiente pede mais uma informação obrigatória que é onde nós vamos dizer quantas colunas vai possuir nossa tabela e juntamente definir a orientação do texto na tabela utilizando os valores "l" de left, "c" de center e "r" de right, e cada uma desses valores vão representar uma coluna como por exemplo \textbackslash begin\{tabular\}\{l c r\} (este comando diz que a tabela vai ter 3 colunas, a primeira com o texto orientado para a esquerda, a segunda com o texto centralizado e a terceira com o texto orientado para a direita), mas além desses valores nós podemos desenhar divisórias entre as colunas com a barra reta "|" como por exemplo \textbackslash begin\{tabular\}\{|l|c|r|\}, e linhas utilizando "\textbackslash hline". Para escrevermos os valores dentro dessas colunas e para dividirmos as linhas nós utilizamos dois valores, o "\&" ele é usado para dividir os conteúdos das colunas e o "\textbackslash \textbackslash" é utilizado para quebrar a linha. este é um exemplo de tabela com essas informações:

\begin{table}[htbp]
    \centering
    \begin{tabular}{|l|c|r|}
        \hline 
         Título 1 & Título 2 & Título 3 \\
         \hline
         % Conteúdo orientado para esquerda & Conteúdo centralizado & Conteúdo orientado para a direita \\
         Conteúdo à esquerda & Conteúdo centralizado & Conteúdo à direita\\
         \hline
    \end{tabular}
    \caption{Tabela de exemplo}
    \label{tab:tabelaExemplo}
\end{table}

Porém esse estilo de tabela é muito pouco utilizado em publicações de artigos por ter um padrão considerado feio, com espaçamentos irregulares e alguns outros defeitos. Para criar uma tabela com padrão de publicação o que fazemos é primeiramente importar um pacote chamado booktabs (\textbackslash usepackage\{booktabs\}). Este pacote permite a gente usar 3 funções, o \textbackslash toprule, \textbackslash midrule e o \textbackslash bottomrule, então para a tabela ficar em padrão de publicação o recomendado a se fazer é utilizar o \textbackslash toprule antes dos títulos da tabela, o \textbackslash midrule logo abaixo do título da tabela e o \textbackslash bottomrule no final da tabela, e por fim esquecer a utilização de linhas para dividir as colunas, ficando mais ou menos assim:

\begin{table}[htbp]
    \centering
    \caption{Este é um exemplo de tabela com padrão de publicação.}
    \begin{tabular}{c c}
        \toprule
        Título 1 & Título 2 \\
        \midrule
         Conteúdo 1.1 & Conteúdo 2.1\\
         Conteúdo 1.2 & Conteúdo 2.2\\
         Conteúdo 1.3 & Conteúdo 2.3\\
         \bottomrule
    \end{tabular}
    \label{tab:placeholder}
\end{table}

\section{Aula 6}
A aula 6 foi focado nas referências cruzadas e bibliografia

Para fazer as referências cruzadas nó utilizamos o \textbackslash label\{\} que traduzindo seria etiqueta, dentro das chaves nós vamos passar o nome que a gente quer etiquetar a nossa imagem/seção/tabela/figura, comumente nós usamos prefixos para agrupar essas etiquetas, os prefixos mais comuns são "fig:" para figuras e imagens, "tab:" é para tabelas, "eq:" é para equações e "sec:" é para seções. Dessa forma as figuras vão ter suas próprias numerações sequenciais separadas das tabelas, equações e seções, evitando que por exemplo a primeira imagem do documento começe como 3 por conta de duas tabelas que vieram antes.

Para que essa referenciação desses elementos sejam um hiperlink onde quando a gente clicar na referenciação a gente seja mandado diretamente para o próprio elemento nós precisaremos importar um pacote chamado hyperref, e assim que importarmos esse pacote as referencias cruzadas já automaticamente vão se tornar hiperlinks, por exemplo, citando a imagem do professor figura~\ref{fig:zickinho} vocês podem clicar e vão ser redirecionados para a imagem.

Outra coisa que aprendemos nessa aula é a bibliografia e citação de outros trabalhos, o que é muito complicado e cheio de regras no Word porém no \LaTeX isso se torna algo muito simples e automático. O \LaTeX utiliza um arquivo separado do documento principal que é o bibtex, neste bibtex nós vamos adicionando trabalhos que queremos citar e referenciar.

Como Daniel disse~\cite{michels2025real}
\section{Aula 7}
Nesta aula aprendemos sobre criação e utilização de classes documentais. Nas outras aulas nós já aprendemos sobre alguns modelos de documentos, como o book, article, report e beamer, esses modelos possuem um padrão de espaçamento, margens, tamanho de texto e comportamentos únicos, como por exemplo o book, onde ele por padrão só cria capítulos novos em páginas da direita, ou seja páginas pares, e esses modelos são criados a partir de uma classe documental.

Porém cada conferência, universidade, revista, livros, e outros meios de publicação ou até países possuem alguma variação de modelo, template, padrões e regras, por isso elas criam esses templates em \LaTeX  para que a gente não precise ficar ajustando cada detalhe na mão, assim a gente só adicione o texto e as imagens que a gente quer e esses textos e imagens já vão automaticamente seguir todas as regras e os padrões exigidos.

Esses templates são criados a partir de três arquivos, o .cls, o .sty e o .bst. O arquivo .cls é o que vai padronizar a classe documental, é nele que a gente cria as regras de estrutura, layout e formatação de cada elemento do texto, como margens, cabeçalhos, rodapés, tamanho de fontes e etc. O arquivo .sty serve pra importar pacotes de estilos adicionais que são os arquivos que a gente adicionava no preâmbulo, dessa forma a gente separa as importações e deixa o documento mais limpo. Por fim o artigo .bst serve para definir o estilo da bibliografia.

Para utilizar esses arquivos de modelagem de templates é muito simples, logo na primeira linha do nosso documento principal onde a gente define o nosso \textbackslash documentclass nós passamos como argumento o nome do nosso arquivo .cls porém sem a extensão .cls, por exemplo, se o arquivo se chama template-udesc.cls nós importamos ele com \textbackslash documentclass\{template-udesc\}. Para adicionar nossos estilos nós utilizamos o próprio \textbackslash usepackage\{\} passando como argumento somente o nome do nosso arquivo .sty sem a extensão. E por fim para usarmos a estilização das nossas referências bibliográficas nós utilizamos o \textbackslash bibliographystyle\{\} passando como arguemento o nome do arquivo .bst sem extensão e depois utilizamos o \textbackslash bibliography\{\} com o nome do nosso arquivo .bib.



\section{Aula 8}
Na aula 8 a gente começa a sair do overleaf e começamos a aprender a instalar o \LaTeX na própria máquina, mas pra que a gente faria isso sendo que podemos usar o overleaf que é muito mais prático na questão de ser online e a gente poder trabalhar em documentos de qualquer lugar? A resposta disso é que o overleaf não é de graça e existem limitações no plano gratis do overleaf como o problema de timeout de compilação, onde dependendo do tamanho e peso do seu documento o overleaf não compila o documento, tem o problema em questão de privacidade e sigilo onde os dados sensíveis de documentos vão ter que trafegar em servidores de terceiros e por fim temos o problema de falta de integração externa onde com o \LaTeX local nos permite fazer integração por exemplo com o python para gerar tabelas dinâmicas direto no \LaTeX.

Primeiramente a gente tem que ter em mente que o \LaTeX não é um programa em si, a gente não vai abrir uma interface pra conseguir desenvolver documentos em \LaTeX, o que acontece é que o \LaTeX é basicamente uma linguagem de programação como o python por exemplo, a gente não baixa o python e abre o python, o python fica em "segundo plano", é apenas um backend, um compilador que a gente não vê de fato.

Existem 3 principais motores de compilação do \LaTeX, que vão ser como o "cérebro" que vai ler o código e transformar ele em pdf, os três motores são o pdfLaTeX que é o mais antigo e mais estável, porém por ser antigo tem alguns problemas tipo por exemplo ele não possui a fonte arial nativamente o que é um grande problema visto que atualmente a grande maioria dos documentos exige fonte arial, o segndo compilador é o XeLaTeX que é considerado um motor mais moderno que já possui a fonte arial e suporte UTF-8, e por fim temos o terceiro compilador que é o LuaLaTeX considerado o futuro do \LaTeX, é um motor embutido com a liguagem de programação Lua permitindo cálculos absurdos durante a compilação.

Ao baixar o \LaTeX a gente pode usar algumas IDEs que podemos usar para editar os documentos que são específicas para utilizar o \LaTeX, como por exemplo o TeXstudio, mas nós podemos usar também o editor vscode, e para isso nós precisamos ter o \LaTeX instalado na máquina e precisamos também instalar a extensão chamada LaTeX WorkShop que vai trazer várias funcionalidades para gente, como o Auto-build, onde logo quando salvarmos o arquivo o nosso pdf já compila e atualiza automaticamente, o SyncTeX que conecta o pdf ao código onde quando damos um duplo clique no pdf ele já joga a gente para a linha exata do código, e também outro exemplo é o linter que vai sublinhar erros de digitação e formatação do nosso arquivo.

Rodando o \LaTeX localmente a gente vai ter que se acostumar a ler os logs de problemas no "modo raiz", esses erros vão ser descritos no arquivo .log ou nós podemos ver eles usando o output da compilação. Alguns dos erros mais comuns que podem aparecer são o "Overfull \textbackslash hbox (15.23pt too wide) in paragraph at lines 45--48" que indica que algum elemento do texto, seja ele uma figura, tabela ou até mesmo alguma palavra extrapolou as margens invisíveis de recuo do documento, outro erro é o "Missing \$ inserted." que indica que a gente usou algum tipo de símbolo matemático no meio do texto sem abrir aspas matemáticas (\$), outro é o "Undefined control sequence." que indica que digitamos algum comando errado e por fim o "File ended while scanning use of..." que indica que teve alguma ambiente aberto que não foi fechada.


\section{Aula 9}
Nesta aula nós saímos um pouco do \LaTeX e começamos a aprender sobre Git e GitHub

\section{Aula 10}
\lipsum[1 - 3]

\section{Aula 11}
\lipsum[1 - 3]

\section{Aula 12}
\lipsum[1 - 3]
